\subsection{Evolution of reduced-order solutions without conserved quantities}

Consider the instantaneous error:
\begin{equation}
    \label{eq:J1}
    \mathcal{J} (\mathbf{q}, \dot{\mathbf{q}}) = \frac{1}{2} \|\hat{u}_t - F(\hat{u})\|^2_H
\end{equation} 

\begin{lemma}[Metric Tensor]
    \label{lemma:metric_tensor}
    Let \autoref{assum:1} hold. Then, the metric tensor $M$ is defined as
    \begin{equation}
        M \in \mathbb{R}^{n\times n} :\hspace{20pt} M_{ij} = \Biggl \langle\pd{\hat{u}}{q_i}, \pd{\hat{u}}{q_j} \Biggr \rangle _H
        \hspace{20pt} \forall i,j \in \{1, 2, ... , n\}
    \end{equation}
    is a symmetric positive-definite matrix for all $\mathbf{q}\in\Omega$.
\end{lemma}
% \label{lemma:metric_tensor}

\begin{proof}
    From definition it is trivial that the matrix $M$ is symmetric. 
    To prove that $M$ is positive-definite:
    \begin{equation}
        \langle \boldsymbol{\xi}, M \boldsymbol{\xi} \rangle_H = 
        \biggl \langle \pd{\hat{u}}{q_i}\xi_{i},\pd{\hat{u}}{q_j}\xi_{j} \biggr \rangle_H = 
        \biggl \| \pd{\hat{u}}{q_i}\xi_{i} \biggr \|_H^2   > 0 
    \end{equation}
\end{proof}

\begin{equation}
    \min_{\dot{\mathbf q} \in \mathbb{R}^n} \mathcal{J}(\mathbf q, \dot{\mathbf q})
    \label{eq:min_pb1}
\end{equation}

\begin{theorem}
    Assuming that \autoref{assum:1} is satisfied, the minimization problem defined in \eqref{eq:min_pb1} admits a unique solution. 
    Moreover, the optimal parameter evolution is given by:
    
    \begin{equation}
        \dot{\mathbf q}=M^{-1}(\mathbf q),\mathbf f(\mathbf q),
        \label{eq:parameter_evolution}
    \end{equation}
    
    where $$M(\mathbf q)$$ denotes the metric tensor introduced in \autoref{lemma:metric_tensor}, and $\mathbf f:\mathbb R^n\rightarrow\mathbb R^n$ is defined component-wise as
    
    \begin{equation}
        f_i = \biggl \langle 
        \frac{\partial \hat u}{\partial q_i},
        F(\hat u)
        \biggl \rangle _H,
        \qquad i=1,\ldots,n.
    \end{equation}
    
\end{theorem}
    
\begin{proof}
    The instantaneous error functional introduced in \eqref{eq:J1} can be expanded as

    \begin{align*}
        J(\mathbf q,\dot{\mathbf q})
        &=
        \frac12 \Big[
        \langle \hat u_t,\hat u_t\rangle_H
        -2\langle \hat u_t, F(\hat u)\rangle_H
        +\langle F(\hat u), F(\hat u)\rangle_H
        \Big] \\
        &=
        \frac12 \dot q_i \dot q_j
        \left\langle
        \frac{\partial \hat u}{\partial q_i},
        \frac{\partial \hat u}{\partial q_j}
        \right\rangle_H
        -
        \dot q_i
        \left\langle
        \frac{\partial \hat u}{\partial q_i},
        F(\hat u)
        \right\rangle_H
        +
        \frac12 \langle F(\hat u), F(\hat u)\rangle_H \\
        &=
        \frac12 \left\langle \dot{\mathbf q}, M(\mathbf q)\dot{\mathbf q} \right\rangle
        -
        \left\langle \dot{\mathbf q}, \mathbf f(\mathbf q) \right\rangle
        +
        \frac12 \langle F(\hat u), F(\hat u)\rangle_H.
    \end{align*}

    Therefore, $\mathcal J$ is a smooth quadratic function of $\dot{\mathbf q}$. 
    Since the metric tensor $M(\mathbf q)$ is symmetric and positive definite, the functional is strictly convex with respect to $\dot{\mathbf q}$. 
    Consequently, the minimization problem possesses a unique solution, characterized by the first-order optimality condition:
    
    \begin{equation}
    \nabla_{\dot{\mathbf q}} J(\mathbf q,\dot{\mathbf q}) = 0.
    \end{equation}
    
    This condition yields
    
    \begin{equation}
    M(\mathbf q)\dot{\mathbf q}-\mathbf f(\mathbf q)=0.
    \end{equation}
    
    Finally, because $M(\mathbf q)$ is invertible, the parameter dynamics can be written as
    
    \begin{equation}
    \dot{\mathbf q}=M^{-1}(\mathbf q)\mathbf f(\mathbf q),
    \end{equation}
    
    which proves the result.
\end{proof}




\subsection{Evolution of reduced-order solutions with conserved quantities}